% GENERATED FILE. Edit canonical lessons, then run npm run generate:books.
% canonical-source-hash: fnv1a64:ff79ce557b55f4c7
% canonical-lessons: FA-C251-honarmand, FA-C251-davar, FA-C251-morabbi, FA-C251-bazikon, FA-C251-tarikh

\chapter{Artist, Referee, Coach, Player, History}
\label{ch:fa-artist-referee-coach-player-history}
\hlchaptermodality{251}

\begin{chapteropening}
\textbf{By the end of this chapter:} \emph{I can say an artist, a referee, a coach, a player and history.}

Complete the last lesson of chapter 251: I can say an artist, a referee, a coach, a player and history.
\end{chapteropening}

\section[\texorpdfstring{\fa{هنرمند} (honarmand)}{هنرمند (honarmand)}]{\fa{هنرمند} (honarmand) — an artist}

\label{lesson:FA-C251-honarmand}



\begin{quote}
Before the new one: say the Persian for a university lecturer, then the Persian for a pilot.
\end{quote}

\subsection*{You'll want to know: \fa{هنرمند}}
\textbf{\fa{هنرمند}} — \emph{honarmand} — \textquotedblleft{}an artist\textquotedblright{}.

\begin{cousinweb}
\emph{honarmand} is \emph{honar}, \textquotedblleft{}art, skill\textquotedblright{}, with the ending \emph{-mand}, \textquotedblleft{}having\textquotedblright{}: one who has art. \emph{honar} is native, usually traced to Middle Persian \emph{hunar}, \textquotedblleft{}skill, virtue\textquotedblright{}.
\end{cousinweb}

\begin{grammarlens}[title={What you've built}]
One more word for this chapter.
\end{grammarlens}

\subsection*{Your turn}
\begin{itemize}
\raggedright
  \item \emph{Say it:} \emph{honarmand}
  \item \emph{Say it:} \emph{honarmand}, once more
  \item \emph{Say it:} say \emph{honarmand}
  \item \emph{Recall:} say the Persian for a university lecturer, then the Persian for a pilot, then say \emph{honarmand} again
\end{itemize}

\begin{tcolorbox}[breakable,colback=teal!4,colframe=teal!35!black,title={Before you move on}]
What does \fa{هنرمند} mean? (\textquotedblleft{}An artist\textquotedblright{}.) Say it once more.
\end{tcolorbox}
\section[\texorpdfstring{\fa{داور} (dâvar)}{داور (dâvar)}]{\fa{داور} (dâvar) — a referee}

\label{lesson:FA-C251-davar}



\begin{quote}
Before the new one: say the Persian for a pilot, then the Persian for an artist.
\end{quote}

\subsection*{You'll want to know: \fa{داور}}
\textbf{\fa{داور}} — \emph{dâvar} — \textquotedblleft{}a referee\textquotedblright{}.

\begin{cousinweb}
\emph{dâvar} is a native Persian word, inherited rather than borrowed; it is usually traced to Middle Persian \emph{dādwar}, \textquotedblleft{}judge\textquotedblright{}.
\end{cousinweb}

\begin{grammarlens}[title={What you've built}]
One more word for this chapter.
\end{grammarlens}

\subsection*{Your turn}
\begin{itemize}
\raggedright
  \item \emph{Say it:} \emph{dâvar}
  \item \emph{Say it:} \emph{dâvar}, once more
  \item \emph{Say it:} say \emph{dâvar}
  \item \emph{Recall:} say the Persian for a pilot, then the Persian for an artist, then say \emph{dâvar} again
\end{itemize}

\begin{tcolorbox}[breakable,colback=teal!4,colframe=teal!35!black,title={Before you move on}]
What does \fa{داور} mean? (\textquotedblleft{}A referee\textquotedblright{}.) Say it once more.
\end{tcolorbox}
\section[\texorpdfstring{\fa{مربی} (morabbi)}{مربی (morabbi)}]{\fa{مربی} (morabbi) — a coach, a trainer}

\label{lesson:FA-C251-morabbi}



\begin{quote}
Before the new one: say the Persian for an artist, then the Persian for a referee.
\end{quote}

\subsection*{You'll want to know: \fa{مربی}}
\textbf{\fa{مربی}} — \emph{morabbi} — \textquotedblleft{}a coach, a trainer\textquotedblright{}.

\begin{cousinweb}
\emph{morabbi} is borrowed from Arabic \emph{murabbī}, \textquotedblleft{}one who brings up, educator\textquotedblright{}. Its Arabic root is \emph{r-b-w} (Lane).
\end{cousinweb}

\begin{grammarlens}[title={What you've built}]
One more word for this chapter.
\end{grammarlens}

\subsection*{Your turn}
\begin{itemize}
\raggedright
  \item \emph{Say it:} \emph{morabbi}
  \item \emph{Say it:} \emph{morabbi}, once more
  \item \emph{Say it:} say \emph{morabbi}
  \item \emph{Recall:} say the Persian for an artist, then the Persian for a referee, then say \emph{morabbi} again
\end{itemize}

\begin{tcolorbox}[breakable,colback=teal!4,colframe=teal!35!black,title={Before you move on}]
What does \fa{مربی} mean? (\textquotedblleft{}A coach, a trainer\textquotedblright{}.) Say it once more.
\end{tcolorbox}
\section[\texorpdfstring{\fa{بازیکن} (bâzikon)}{بازیکن (bâzikon)}]{\fa{بازیکن} (bâzikon) — a player}

\label{lesson:FA-C251-bazikon}



\begin{quote}
Before the new one: say the Persian for a referee, then the Persian for a coach.
\end{quote}

\subsection*{You'll want to know: \fa{بازیکن}}
\textbf{\fa{بازیکن}} — \emph{bâzikon} — \textquotedblleft{}a player\textquotedblright{}.

\begin{cousinweb}
\emph{bâzikon} is built from \emph{bâzi}, \textquotedblleft{}play\textquotedblright{}, and \emph{-kon}, \textquotedblleft{}doing\textquotedblright{} (from \emph{kardan}): the one who plays. \emph{bâzi} is usually traced to Middle Persian \emph{wāzīg}, and \emph{kardan} to \emph{kardan} (EtymDB).
\end{cousinweb}

\begin{grammarlens}[title={What you've built}]
One more word for this chapter.
\end{grammarlens}

\subsection*{Your turn}
\begin{itemize}
\raggedright
  \item \emph{Say it:} \emph{bâzikon}
  \item \emph{Say it:} \emph{bâzikon}, once more
  \item \emph{Say it:} say \emph{bâzikon}
  \item \emph{Recall:} say the Persian for a referee, then the Persian for a coach, then say \emph{bâzikon} again
\end{itemize}

\begin{tcolorbox}[breakable,colback=teal!4,colframe=teal!35!black,title={Before you move on}]
What does \fa{بازیکن} mean? (\textquotedblleft{}A player\textquotedblright{}.) Say it once more.
\end{tcolorbox}
\section[\texorpdfstring{\fa{تاریخ} (târikh)}{تاریخ (târikh)}]{\fa{تاریخ} (târikh) — history; a date}

\label{lesson:FA-C251-tarikh}



\begin{quote}
Before the new one: say the Persian for a coach, then the Persian for a player.
\end{quote}

\subsection*{You'll want to know: \fa{تاریخ}}
\textbf{\fa{تاریخ}} — \emph{târikh} — \textquotedblleft{}history; a date\textquotedblright{}.

\begin{cousinweb}
\emph{târikh} is borrowed from Arabic \emph{tārīḫ}, \textquotedblleft{}date, history\textquotedblright{}. Its Arabic root is \emph{ʾ-r-ḫ} (Lane).
\end{cousinweb}

\begin{grammarlens}[title={What you've built}]
One more word for this chapter.
\end{grammarlens}

\subsection*{Your turn}
\begin{itemize}
\raggedright
  \item \emph{Say it:} \emph{târikh}
  \item \emph{Say it:} \emph{târikh}, once more
  \item \emph{Say it:} say \emph{târikh}
  \item \emph{Recall:} say the Persian for a coach, then the Persian for a player, then say \emph{târikh} again
\end{itemize}

\begin{tcolorbox}[breakable,colback=teal!4,colframe=teal!35!black,title={Before you move on}]
What does \fa{تاریخ} mean? (\textquotedblleft{}History; a date\textquotedblright{}.) Say it once more.
\end{tcolorbox}
